\begin{pastexam}[label={px:ntu-109-comm-dsa-15}]{臺大電信所 109 學年度（B）第 15 題}
（單選；每題 2 分，答錯不倒扣）

In questions that ask for time complexities, please choose the tightest bound.

In a binomial heap of size n, the worst case time complexity of one delete-min operation is

(A)~$O(1)$ \qquad (B)~$O(\log^* n)$ \qquad (C)~$O(\log n)$ \qquad (D)~$O(n)$ \qquad (E)~$O(n\log n)$
\end{pastexam}

\begin{pastsoln}
答案是 (C)。

$n$ 個元素的二項堆積最多 $\lfloor \lg n \rfloor + 1$ 棵樹，每棵的階數不超過
$\lfloor \lg n \rfloor$（資結書第~8.5~節）。 delete-min 分三步：
\begin{enumerate}
\item 掃過根串列找最小的根：$O(\log n)$。
\item 拿掉這個根，它的子樹（最多 $\lfloor \lg n \rfloor$ 棵）反過來排成一個二項堆積：
$O(\log n)$。
\item 和剩下的堆積合併，像二進位加法，每個階數做常數的工作：$O(\log n)$。
\end{enumerate}
合計 $O(\log n)$。 $n = 2^m - 1$ 時有 $m$ 棵樹，光掃根串列就要 $\Theta(\log n)$，所以這個界是緊的；
(A)、(B) 太小，(D)、(E) 不緊。
\end{pastsoln}
